\documentclass[12pt]{article}
\usepackage{amsmath,amsthm,amssymb,amsfonts,setspace}
\usepackage{fancyhdr}
\usepackage{hyperref}

\textwidth 7.0 truein
\oddsidemargin -0.25in
\evensidemargin -0.5 truein
\topmargin -0.85in
\textheight 9.5in

%%%% Custom commands %%%%

% Box/square the final answer
\newcommand{\ans}[1]{\boxed{#1}}

% Shared indentation lengths for subproblems and solutions
\newlength{\subproblemindent}
\setlength{\subproblemindent}{1.5em}
\newlength{\subproblemlabelwidth}
\setlength{\subproblemlabelwidth}{2em}
\newlength{\solutionindent}
\setlength{\solutionindent}{3em}

% Make \circ usable in ordinary text as well as math mode
\let\latexcirc\circ
\renewcommand{\circ}{\ifmmode\latexcirc\else\ensuremath{\latexcirc}\fi}

% Bold a top-level problem: number and text separated by an em-dash
\newcommand{\problem}[2]{\vspace{1em}\noindent\textbf{#1} --- #2\par\vspace{0.25em}}

% Bold and indent a subproblem: letter and text separated by an em-dash
\newcommand{\subproblem}[2]{%
  \par\vspace{0.5em}%
  \noindent
  \hangindent=\dimexpr\subproblemindent+\subproblemlabelwidth+0.5em\relax
  \hangafter=1
  \hspace*{\subproblemindent}%
  \makebox[\subproblemlabelwidth][l]{\textbf{#1}}#2\par
}

% Indent all lines of the answer/solution block
\newcommand{\solution}[1]{%
  \par\vspace{0.0em}%
  \begin{list}{}{
    \setlength{\leftmargin}{\solutionindent}
    \setlength{\rightmargin}{0pt}
    \setlength{\listparindent}{0pt}
    \setlength{\itemindent}{0pt}
    \setlength{\parsep}{0pt}
  }%
    \item\relax
    \begingroup
      \setlength{\baselineskip}{1.5\baselineskip}
      \setlength{\lineskip}{1.5\lineskip}
      \setlength{\parskip}{0pt}
      \everymath{\displaystyle}
      #1
      \par
    \endgroup
  \end{list}%
  \vspace{0.25em}%
}

% Running footer with a link back to the source repo, shown on every page
\pagestyle{fancy}
\fancyhf{}
\renewcommand{\headrulewidth}{0pt}
\renewcommand{\footrulewidth}{0.4pt}
\fancyfoot[C]{\footnotesize\textit{Source: \url{https://github.com/chrismcarrico/homework/blob/master/spce5085/a4/carrico_spce5085_a4.ipynb}}}

%%%% In most cases you won't need to edit anything above this line %%%%

\begin{document}
\hfill Christian Carrico

\hfill SPCE 5085

\hfill Assignment \#4

\problem{7.2}{Why is attenuation on a slant path in clear sky conditions greater at 22 GHz than at 30 GHz? Why is the occurrence of attenuation in rain on a slant path at the 1 - 2 dB level often greater at 22 GHz than at 30 GHz?}
\solution{}

\problem{7.4}{An earth station at sea level communicates at an elevation angle of 30\circ with a GEO satellite in Ka-band at a downlink frequency of 20 GHz. Stratiform rain is present in the slant path with a melting level height of 2.5 km. Find}
\subproblem{a}{The path length through the rain.}
\solution{}
\subproblem{b}{The path attenuation for a specific attenuation of 1.5 dB/km.}
\solution{}
\problem{3}{Create a plot similar to figure 3.1 Doppler vs SV altitude in the article "The Mathematical Model of Doppler Frequency Shift in Leo At Ku, K and Ka Frequency Bands" (posted in the module 2), for satellite signal transmission from spacecraft to stationary ground station - L band (1.25 GHz), X band (11.0 GHz), and Ku band (16.5 GHz) for SV altitudes of (at a minimum) 1000, 10,000 and 20,000 km. (40 pts)}
\solution{}
\end{document}
